\documentclass[11pt,a4paper]{article}

\usepackage[utf8]{inputenc}
\usepackage[T1]{fontenc}
\usepackage[british]{babel}
\usepackage{lmodern}
\usepackage[margin=2.5cm]{geometry}
\usepackage{amsmath,amssymb}
\usepackage{booktabs}
\usepackage{graphicx}
\usepackage{caption}
\usepackage{listings}
\usepackage{xcolor}
\usepackage{fancyhdr}
\usepackage{microtype}
\usepackage[colorlinks=true,linkcolor=blue!50!black,urlcolor=blue!50!black,citecolor=blue!50!black]{hyperref}

% ---------------------------------------------------------------- metadata ---
\newcommand{\authorname}{Your Name}
\newcommand{\repourl}{https://github.com/<username>/amzn-volatility-risk}

% ----------------------------------------------------------------- layout ----
\setlength{\headheight}{14pt}
\pagestyle{fancy}
\fancyhf{}
\fancyhead[L]{\small AMZN returns: volatility modelling and risk}
\fancyhead[R]{\small\authorname}
\fancyfoot[C]{\thepage}
\captionsetup{font=small,labelfont=bf}

\lstset{
  language=R,
  basicstyle=\ttfamily\footnotesize,
  keywordstyle=\color{blue!60!black},
  commentstyle=\color{green!40!black}\itshape,
  stringstyle=\color{red!50!black},
  breaklines=true,
  frame=single,
  columns=fullflexible,
  showstringspaces=false
}

% Include a generated figure if it exists; otherwise show a placeholder box.
\newcommand{\genfigure}[4][0.95]{%
  \begin{figure}[htbp]\centering
  \IfFileExists{figures/#2.png}%
    {\includegraphics[width=#1\linewidth]{figures/#2.png}}%
    {\fbox{\parbox{0.9\linewidth}{\centering\small\vspace{1.5cm}%
      Figure \texttt{\detokenize{#2}} not found. Run \texttt{make analysis} to generate it.%
      \vspace{1.5cm}}}}%
  \caption{#3}\label{#4}
  \end{figure}}

\title{\textbf{Econometric Time-Series Analysis and Quantitative\\Risk Management of Amazon Daily Returns\\(AMZN, 2019--2021)}}
\author{\authorname}
\date{November 2026}

\begin{document}
\maketitle

\begin{abstract}
\noindent This report applies a standard toolkit of financial econometrics to the daily returns of Amazon (AMZN) between 2019 and 2021. It documents the statistical properties of the returns, identifies a conditional-mean model (ARMA), compares it out of sample with a non-linear neural-network autoregression (NNAR), models conditional volatility with an ARMA(1,1)--GARCH(1,1) specification under skewed Student-$t$ innovations, and quantifies market risk through Value at Risk (VaR) and Expected Shortfall (ES), both over the full sample and in a rolling window. The returns are leptokurtic and display strong volatility clustering; the GARCH model removes the conditional heteroscedasticity and shows highly persistent volatility; and Gaussian risk measures understate tail risk relative to historical and Cornish--Fisher estimates. All computations are reproducible with the \textsf{R} code in the accompanying repository.
\end{abstract}

\tableofcontents
\newpage

%====================================================================
\section{Introduction}
%====================================================================
Financial returns rarely behave like independent, normally distributed shocks. They display fat tails, volatility that changes over time, and extreme events that occur far more often than a Gaussian model would suggest. Ignoring these features leads to misjudged risk.

This project demonstrates a complete and coherent workflow for analysing a single asset. Its objectives are to:
\begin{enumerate}
  \item describe the statistical moments of the return series and test for normality;
  \item test for stationarity and identify a model for the conditional mean;
  \item compare a linear (ARMA) and a non-linear (NNAR) forecasting model out of sample;
  \item test for and model conditional heteroscedasticity with a GARCH specification;
  \item measure market risk with VaR and ES, statically and dynamically.
\end{enumerate}

\paragraph{Data.} The data are the daily adjusted closing prices of AMZN from 1 January 2019 to 31 December 2021, obtained from Yahoo Finance through the \texttt{quantmod} package. Log returns are computed as
\begin{equation}
  r_t = \ln\!\left(\frac{P_t}{P_{t-1}}\right).
\end{equation}
The sample is split chronologically: the first 80\% of the observations form the training sample and the remaining 20\% the test sample. The period includes the COVID-19 market shock of early 2020.

%====================================================================
\section{Theoretical Background}
%====================================================================

\subsection{Moments and normality}
The mean, variance, skewness and kurtosis summarise the shape of a distribution. Under normality the skewness is $0$ and the kurtosis is $3$ (excess kurtosis $0$). Financial returns typically exhibit \emph{leptokurtosis}, i.e.\ excess kurtosis above zero and heavy tails. The Jarque--Bera test \cite{jarque1980} examines jointly whether skewness and excess kurtosis are zero.

\subsection{Stationarity and ARMA models}
A series is weakly stationary if its mean, variance and autocovariances do not depend on time. The Augmented Dickey--Fuller (ADF) test \cite{dickey1979} takes a unit root as its null hypothesis, so rejecting it supports stationarity. An $\mathrm{ARMA}(p,q)$ model is
\begin{equation}
  r_t = c + \sum_{i=1}^{p}\phi_i r_{t-i} + \sum_{j=1}^{q}\theta_j \varepsilon_{t-j} + \varepsilon_t .
\end{equation}
Candidate orders can be explored with the Extended Autocorrelation Function (EACF) and compared with information criteria such as the AIC.

\subsection{Neural-network autoregression}
An $\mathrm{NNAR}(p,k)$ model, as implemented by \texttt{nnetar()} in the \texttt{forecast} package \cite{hyndman2021}, feeds $p$ lagged values into a feed-forward network with one hidden layer of $k$ nodes, so it can capture non-linear dependence that ARMA models cannot. Out-of-sample accuracy is assessed with the root mean squared error,
\begin{equation}
  \mathrm{RMSE}=\sqrt{\frac{1}{n}\sum_{t=1}^{n}\left(r_t-\hat r_t\right)^2},
\end{equation}
which is preferred to the MAPE here because returns are frequently close to zero, making percentage errors unstable.

\subsection{ARCH and GARCH models}
Returns are often serially uncorrelated while their squares are not: calm and turbulent periods cluster together (\emph{volatility clustering}). Engle's ARCH-LM test \cite{engle1982} detects this by regressing squared residuals on their own lags; the statistic $nR^2$ is asymptotically $\chi^2_q$ under the null of no ARCH effects. With an ARMA(1,1) mean, the model used here is
\begin{align}
  r_t &= \mu + \phi\, r_{t-1} + \theta\, \varepsilon_{t-1} + \varepsilon_t, \qquad \varepsilon_t = \sigma_t z_t, \\
  \sigma_t^2 &= \omega + \alpha_1 \varepsilon_{t-1}^2 + \beta_1 \sigma_{t-1}^2 ,
\end{align}
the GARCH(1,1) recursion of Bollerslev \cite{bollerslev1986}. The process is covariance-stationary if $\alpha_1+\beta_1<1$, and the closer the sum is to one, the more persistent volatility shocks are; the half-life of a shock is $\ln(0.5)/\ln(\alpha_1+\beta_1)$ days. A GARCH(1,1) is far more parsimonious than the high-order ARCH($p$) needed to reproduce the same persistence. The standardised innovations $z_t$ follow a skewed Student-$t$ (\texttt{sstd}) distribution in the parametrisation of Fern\'andez and Steel \cite{fernandez1998}, with a shape parameter $\nu$ (degrees of freedom) governing the tails and a skew parameter $\xi$, where $\xi=1$ is the symmetric case.

\subsection{Value at Risk and Expected Shortfall}
Using return units, the $\alpha$-level VaR is the $(1-\alpha)$ quantile of the return distribution, so a loss threshold that is exceeded with probability $1-\alpha$ over the horizon. Expected Shortfall (also CVaR) is the average return conditional on breaching it:
\begin{equation}
  \mathrm{ES}_\alpha = \mathbb{E}\left[\,r_t \mid r_t \le \mathrm{VaR}_\alpha\,\right].
\end{equation}
Three estimators are used. The \emph{Gaussian} estimator assumes normality. The \emph{historical} estimator uses empirical quantiles. The \emph{modified} estimator uses the Cornish--Fisher expansion \cite{favre2002}, which replaces the Gaussian quantile $z=\Phi^{-1}(1-\alpha)$ by
\begin{equation}
  z_{cf} = z + \frac{z^2-1}{6}S + \frac{z^3-3z}{24}K - \frac{2z^3-5z}{36}S^2,
\end{equation}
where $S$ and $K$ are the sample skewness and excess kurtosis, and sets $\mathrm{VaR}=\hat\mu+z_{cf}\hat\sigma$. Unlike VaR, ES is a \emph{coherent} risk measure because it is subadditive \cite{artzner1999,acerbi2002}, and it describes the severity of losses beyond the threshold.

%====================================================================
\section{Methodology}
%====================================================================
The analysis is carried out in \textsf{R} with the packages \texttt{quantmod}, \texttt{tseries}, \texttt{forecast}, \texttt{TSA}, \texttt{rugarch} and \texttt{PerformanceAnalytics}. The workflow is organised as one script per stage (Appendix~\ref{app:repro}):
\begin{enumerate}
  \item \textbf{Data preparation:} download prices, compute log returns, split into training and test samples.
  \item \textbf{Descriptive analysis:} moments, Jarque--Bera test, ADF test, Ljung--Box tests on returns and squared returns, and the EACF.
  \item \textbf{Mean model:} ARMA(1,1) versus MA(1) via the AIC, plus \texttt{auto.arima()}.
  \item \textbf{Non-linear model:} NNAR$(13,7)$ (13 lags, 7 hidden nodes, no seasonality) against the ARMA models, using the out-of-sample RMSE over the whole test sample, with a zero-forecast benchmark.
  \item \textbf{Volatility:} ARCH-LM test on the ARMA(1,1) residuals; estimation of ARMA(1,1)--sGARCH(1,1) with \texttt{sstd} innovations on the training sample; ARCH-LM and Ljung--Box tests on the standardised squared residuals.
  \item \textbf{Risk:} one-day, 95\% VaR and ES with the three estimators over the full sample; rolling-window estimates with a 500-observation window; and an informal exceedance check of next-day returns against the rolling VaR (Kupiec proportion-of-failures test).
\end{enumerate}

%====================================================================
\section{Results}
%====================================================================

\subsection{Descriptive statistics}
The sample mean of the log returns is statistically indistinguishable from zero. The series alternates between quiet spells and bursts of high volatility (Figure~\ref{fig:returns}), a first visual sign of volatility clustering, and the autocorrelation of squared returns is much stronger than that of returns (Figure~\ref{fig:acf}).

\begin{table}[htbp]
\centering
\caption{Descriptive statistics of AMZN daily log returns}\label{tab:desc}
\IfFileExists{tables/descriptive_stats.tex}%
  {\input{tables/descriptive_stats.tex}}%
  {\small\emph{Table not found. Run \texttt{make analysis} to generate it.}}
\end{table}

\genfigure{returns_series}{AMZN daily log returns, 2019--2021.}{fig:returns}
\genfigure{acf_returns}{Sample autocorrelation of returns (left) and squared returns (right).}{fig:acf}

When a skewed Student-$t$ distribution is fitted within the GARCH model (Section~\ref{sec:garch}), the shape parameter is about $5.2$ degrees of freedom, a low value that indicates heavy tails (the fourth moment is finite only for $\nu>4$), while the skew parameter is about $0.97$, close to the symmetric value $\xi=1$, so asymmetry is mild. The Jarque--Bera test rejects normality (Figure~\ref{fig:dist}). The series is therefore \textbf{leptokurtic and non-normal}: extreme observations are more likely than a Gaussian distribution implies.

\genfigure{return_distribution}{Histogram against the fitted normal density (left) and normal Q--Q plot (right).}{fig:dist}

\subsection{Stationarity and mean model}
The ADF test rejects the unit root, so the log returns are stationary, $I(0)$. The EACF suggested an MA(1) or ARMA(1,1) structure, and \texttt{auto.arima()} selected ARMA(1,1), i.e.\ ARIMA$(1,0,1)$. This specification achieved a lower AIC than the fixed MA(1) alternative (Table~\ref{tab:models}).

\subsection{ARMA versus NNAR}
Table~\ref{tab:models} compares the models over the test sample. In the original run, the non-linear NNAR$(13,7)$ model produced the lowest out-of-sample RMSE, which suggests that some non-linear structure is present in the returns. Two caveats apply. First, the forecasts are multi-step over the entire test sample, so the ARMA forecasts quickly converge to the unconditional mean; differences between the models are therefore small in absolute terms and were not tested formally (for instance with a Diebold--Mariano test). Second, the NNAR network is initialised randomly, so its figures depend on the seed (fixed in the code).

\begin{table}[htbp]
\centering
\caption{Mean-model comparison. AIC is computed on the training sample; RMSE and MAE (in percentage points) on the test sample.}\label{tab:models}
\IfFileExists{tables/model_comparison.tex}%
  {\input{tables/model_comparison.tex}}%
  {\small\emph{Table not found. Run \texttt{make analysis} to generate it.}}
\end{table}

\genfigure{forecast_comparison}{Out-of-sample forecasts of the ARMA(1,1) and NNAR models against realised returns.}{fig:fc}

\subsection{Conditional heteroscedasticity and GARCH}\label{sec:garch}
The ARCH-LM test applied to the ARMA(1,1) residuals rejects homoscedasticity at lags 4, 8, 12 and 16 (all $p$-values close to zero), confirming significant ARCH effects. An ARMA(1,1)--sGARCH(1,1) model with \texttt{sstd} innovations was then estimated on the training sample. Table~\ref{tab:garch} summarises the persistence results.

\begin{table}[htbp]
\centering
\caption{Persistence in the estimated sGARCH(1,1) model}\label{tab:garch}
\begin{tabular}{lc}
\toprule
Parameter & Estimate \\
\midrule
$\alpha_1$ (ARCH effect) & 0.1151 \\
$\beta_1$ (GARCH effect) & 0.8695 \\
$\alpha_1+\beta_1$ & 0.9846 \\
\bottomrule
\end{tabular}
\end{table}

Since $\alpha_1+\beta_1=0.9846<1$, the process is stationary and volatility mean-reverts, but the sum is very close to one, which reflects \textbf{high volatility persistence}: a shock to volatility has a half-life of roughly 45 trading days. After estimation, the Ljung--Box test on the squared standardised residuals and the ARCH-LM test both give $p$-values above $0.05$, so the model removes the conditional heteroscedasticity (full diagnostics are written to \texttt{outputs/garch\_diagnostics.csv}). Figure~\ref{fig:vol} shows the estimated conditional volatility.

\genfigure{conditional_volatility}{Estimated conditional standard deviation (training sample) against absolute returns.}{fig:vol}

\subsection{Value at Risk and Expected Shortfall}
Table~\ref{tab:var} shows the one-day, 95\% risk measures over the full sample, expressed as returns so that negative values are losses.

\begin{table}[htbp]
\centering
\caption{One-day 95\% VaR and ES for AMZN log returns (full sample)}\label{tab:var}
\begin{tabular}{lcc}
\toprule
Method & VaR$_{95\%}$ & ES$_{95\%}$ \\
\midrule
Gaussian & $-3.06\%$ & $-3.87\%$ \\
Historical & $-2.90\%$ & $-4.21\%$ \\
Modified (Cornish--Fisher) & $-2.93\%$ & $-4.20\%$ \\
\bottomrule
\end{tabular}
\end{table}

The Gaussian ES is noticeably smaller in magnitude than the historical and modified figures, consistent with the heavy tails found earlier: assuming normality understates the average loss on the worst 5\% of days. The modified estimator, which corrects for skewness and kurtosis, lands very close to the historical result. The rolling-window estimates (Figure~\ref{fig:roll}) show that the risk measures are not constant: they rise in turbulent periods and fall in calmer ones. An exceedance check of next-day returns against the rolling VaR, with a Kupiec test, is reported in \texttt{outputs/var\_backtest\_kupiec.csv}.

\genfigure{rolling_risk}{Daily returns with rolling 95\% one-day VaR and ES (500-day window).}{fig:roll}

%====================================================================
\section{Conclusions}
%====================================================================
\begin{itemize}
  \item AMZN log returns are stationary, have a mean close to zero and are clearly leptokurtic, so the normal distribution is not an adequate description.
  \item ARMA(1,1) is the preferred linear specification for the conditional mean according to the AIC, but the non-linear NNAR$(13,7)$ achieved a lower out-of-sample RMSE.
  \item Strong ARCH effects are present. An ARMA(1,1)--sGARCH(1,1) model with skewed Student-$t$ innovations captures them adequately, and the high persistence ($\alpha_1+\beta_1\approx0.985$) is typical of equity volatility.
  \item Gaussian risk measures understate tail risk. Historical and Cornish--Fisher estimates are more realistic, and ES is the more informative and theoretically sounder measure because it is coherent.
\end{itemize}

\paragraph{Limitations.} The study covers a single asset and a short period that includes the COVID-19 shock. The train/test split is a single chronological split and the forecast comparison is multi-step rather than rolling one-step-ahead. The static VaR and ES are computed in sample, and the rolling VaR is based on empirical or moment-based estimators rather than on the GARCH conditional distribution. Backtesting is limited to an unconditional-coverage (Kupiec) check, without a Christoffersen independence test. Natural extensions are asymmetric volatility models (EGARCH, GJR-GARCH), GARCH-based VaR/ES forecasts, rolling one-step-ahead forecast evaluation, and a multi-asset portfolio.

%====================================================================
\appendix
\section{Reproducibility}\label{app:repro}
%====================================================================
The complete code is available at \url{\repourl}. After installing the required packages (\texttt{make install}), the whole pipeline is run with \texttt{make analysis} and this report is compiled with \texttt{make report}. Each script writes its outputs to \texttt{outputs/} (CSV and text), \texttt{report/figures/} and \texttt{report/tables/}; the session information is stored in \texttt{outputs/session\_info.txt}.

\begin{table}[htbp]
\centering
\caption{Scripts and their outputs}
\begin{tabular}{ll}
\toprule
Script & Purpose \\
\midrule
\texttt{R/00\_config.R} & Parameters, data download and caching, helper functions \\
\texttt{R/01\_descriptive.R} & Moments, normality, stationarity, EACF, figures \\
\texttt{R/02\_mean\_models.R} & ARMA, \texttt{auto.arima}, NNAR, forecast comparison \\
\texttt{R/03\_garch.R} & ARCH-LM tests, ARMA--sGARCH estimation, diagnostics \\
\texttt{R/04\_risk.R} & Static and rolling VaR/ES, Kupiec exceedance check \\
\texttt{R/run\_all.R} & Runs the pipeline in order \\
\bottomrule
\end{tabular}
\end{table}

\section{Key code excerpts}
\begin{lstlisting}
# ARMA(1,1)-sGARCH(1,1) with skewed Student-t innovations (R/03_garch.R)
spec <- ugarchspec(
  variance.model     = list(model = "sGARCH", garchOrder = c(1, 1)),
  mean.model         = list(armaOrder = c(1, 1), include.mean = TRUE),
  distribution.model = "sstd"
)
fit <- ugarchfit(spec = spec, data = train)

# Rolling VaR and ES, 500-observation window (R/04_risk.R)
risk_measures <- function(x, p) {
  x <- as.numeric(x)
  c(VaR_gaussian   = as.numeric(VaR(x, p = p, method = "gaussian")),
    VaR_historical = as.numeric(VaR(x, p = p, method = "historical")),
    VaR_modified   = as.numeric(VaR(x, p = p, method = "modified")),
    ES_gaussian    = as.numeric(ES(x,  p = p, method = "gaussian")),
    ES_historical  = as.numeric(ES(x,  p = p, method = "historical")),
    ES_modified    = as.numeric(ES(x,  p = p, method = "modified")))
}
roll <- rollapply(as.zoo(d$ret), width = 500, FUN = risk_measures,
                  p = 0.95, by.column = FALSE, align = "right")
\end{lstlisting}

%====================================================================
\begin{thebibliography}{99}
%====================================================================
\bibitem{acerbi2002} Acerbi, C. and Tasche, D. (2002). On the coherence of expected shortfall. \emph{Journal of Banking \& Finance}, 26(7), 1487--1503.
\bibitem{artzner1999} Artzner, P., Delbaen, F., Eber, J.-M. and Heath, D. (1999). Coherent measures of risk. \emph{Mathematical Finance}, 9(3), 203--228.
\bibitem{bollerslev1986} Bollerslev, T. (1986). Generalized autoregressive conditional heteroskedasticity. \emph{Journal of Econometrics}, 31(3), 307--327.
\bibitem{dickey1979} Dickey, D. A. and Fuller, W. A. (1979). Distribution of the estimators for autoregressive time series with a unit root. \emph{Journal of the American Statistical Association}, 74(366), 427--431.
\bibitem{engle1982} Engle, R. F. (1982). Autoregressive conditional heteroscedasticity with estimates of the variance of United Kingdom inflation. \emph{Econometrica}, 50(4), 987--1007.
\bibitem{favre2002} Favre, L. and Galeano, J.-A. (2002). Mean-modified Value-at-Risk optimization with hedge funds. \emph{Journal of Alternative Investments}, 5(2), 21--25.
\bibitem{fernandez1998} Fern\'andez, C. and Steel, M. F. J. (1998). On Bayesian modeling of fat tails and skewness. \emph{Journal of the American Statistical Association}, 93(441), 359--371.
\bibitem{hyndman2021} Hyndman, R. J. and Athanasopoulos, G. (2021). \emph{Forecasting: Principles and Practice}, 3rd edition. OTexts, Melbourne.
\bibitem{jarque1980} Jarque, C. M. and Bera, A. K. (1980). Efficient tests for normality, homoscedasticity and serial independence of regression residuals. \emph{Economics Letters}, 6(3), 255--259.
\end{thebibliography}

\end{document}
